\documentclass[12pt,a4paper]{article}
\usepackage[utf8]{inputenc}
\usepackage[T2A,T1]{fontenc}
\usepackage[vietnamese,english]{babel}
\usepackage{geometry}
\geometry{a4paper, margin=1in}
\usepackage{listings}
\usepackage{xcolor}
\usepackage{hyperref}

\definecolor{codegreen}{rgb}{0,0.6,0}
\definecolor{codegray}{rgb}{0.5,0.5,0.5}
\definecolor{codepurple}{rgb}{0.58,0,0.82}
\definecolor{backcolour}{rgb}{0.95,0.95,0.92}

\lstdefinestyle{mystyle}{
    backgroundcolor=\color{backcolour},   
    commentstyle=\color{codegreen},
    keywordstyle=\color{magenta},
    numberstyle=\tiny\color{codegray},
    stringstyle=\color{codepurple},
    basicstyle=\ttfamily\footnotesize,
    breakatwhitespace=false,         
    breaklines=true,                 
    captionpos=b,                    
    keepspaces=true,                 
    numbers=left,                    
    numbersep=5pt,                  
    showspaces=false,                
    showstringspaces=false,
    showtabs=false,                  
    tabsize=2
}
\lstset{style=mystyle}

\title{ASIC Verification - Exercise 1 Report}
\author{V\~u Ti\~en Giang \\ MSHV: 2570188}
\date{\today}

\begin{document}

\maketitle

\section{Exercise 1: Counter Design and Verification}

\subsection{Objective}
Design a 5-bit synchronous counter based on the specifications provided and write a testbench to verify its correctness against a target timing trace.

\subsection{Design Specifications}
\begin{itemize}
    \item The counter is clocked on the rising edge of \texttt{clk}.
    \item \texttt{rst} is active high and asynchronously resets the output to 0.
    \item \texttt{cnt\_in} and \texttt{cnt\_out} are 5-bit signals.
    \item If \texttt{load} is high, the counter is loaded from \texttt{cnt\_in} (takes priority over \texttt{enab}).
    \item Otherwise, if \texttt{enab} is high, \texttt{cnt\_out} is incremented.
\end{itemize}

\subsection{Implementation (\texttt{counter.v})}
\begin{lstlisting}[language=Verilog]
`timescale 1ns / 1ps

module counter (
    input clk,
    input rst,
    input load,
    input enab,
    input [4:0] cnt_in,
    output reg [4:0] cnt_out
);
    always @(posedge clk or posedge rst) begin
        if (rst) begin
            cnt_out <= 5'b00000;
        end else if (load) begin
            cnt_out <= cnt_in;
        end else if (enab) begin
            cnt_out <= cnt_out + 1;
        end
    end
endmodule
\end{lstlisting}

\subsection{Testbench (\texttt{counter\_tb.v})}
To match the precise timestamps detailed in the exercise, the testbench drives inputs securely synchronized around the clock edges.

\begin{lstlisting}[language=Verilog]
`timescale 1ns / 1ps

module counter_tb;
    reg clk;
    reg rst, load, enab;
    reg [4:0] cnt_in;
    wire [4:0] cnt_out;

    counter uut (.clk(clk), .rst(rst), .load(load), .enab(enab), .cnt_in(cnt_in), .cnt_out(cnt_out));

    initial begin
        clk = 0;
        forever #5 clk = ~clk;
    end

    initial begin
        rst = 1; load = 0; enab = 0; cnt_in = 0;

        #12; rst = 0; load = 1; enab = 1; cnt_in = 5'b10101; 
        #8;  $display("At time %0t rst=%b load=%b enab=%b cnt_in=%b cnt_out=%b", 
             $time, rst, load, enab, cnt_in, cnt_out);
        
        #2;  rst = 0; load = 1; enab = 1; cnt_in = 5'b01010;
        #8;  $display("At time %0t rst=%b load=%b enab=%b cnt_in=%b cnt_out=%b", 
             $time, rst, load, enab, cnt_in, cnt_out);
        
        #2;  rst = 0; load = 1; enab = 1; cnt_in = 5'b11111;
        #8;  $display("At time %0t rst=%b load=%b enab=%b cnt_in=%b cnt_out=%b", 
             $time, rst, load, enab, cnt_in, cnt_out);
        
        #2;  rst = 1; load = 1; enab = 1; cnt_in = 5'b11111;
        #8;  $display("At time %0t rst=%b load=%b enab=%b cnt_in=%b cnt_out=%b", 
             $time, rst, load, enab, cnt_in, cnt_out);
        
        #2;  rst = 0; load = 1; enab = 1; cnt_in = 5'b11111;
        #8;  $display("At time %0t rst=%b load=%b enab=%b cnt_in=%b cnt_out=%b", 
             $time, rst, load, enab, cnt_in, cnt_out);
        
        #2;  rst = 0; load = 0; enab = 1; cnt_in = 5'b11111;
        #8;  $display("At time %0t rst=%b load=%b enab=%b cnt_in=%b cnt_out=%b", 
             $time, rst, load, enab, cnt_in, cnt_out);

        #10; $finish;
    end
endmodule
\end{lstlisting}

\subsection{Simulation Results}
The simulation log matches the specified requirements at the detailed time points.
\begin{lstlisting}[language=bash]
At time 20000 rst=0 load=1 enab=1 cnt_in=10101 cnt_out=10101
At time 30000 rst=0 load=1 enab=1 cnt_in=01010 cnt_out=01010
At time 40000 rst=0 load=1 enab=1 cnt_in=11111 cnt_out=11111
At time 50000 rst=1 load=1 enab=1 cnt_in=11111 cnt_out=00000
At time 60000 rst=0 load=1 enab=1 cnt_in=11111 cnt_out=11111
At time 70000 rst=0 load=0 enab=1 cnt_in=11111 cnt_out=00000
\end{lstlisting}

\section{Exercise 2: Stack Design and Verification}

\subsection{Objective}
Create a parameterized LIFO (Last-In-First-Out) Stack design, present testing scenarios, and write a testbench to verify edge cases.

\subsection{Design Specifications}
\begin{itemize}
    \item Parameterized for \texttt{WIDTH} and \texttt{DEPTH}. \texttt{DEPTH} is dynamically managed via \texttt{DEPTH\_LOG2} (power of 2 constraint).
    \item Asynchronous high-active reset (\texttt{rst}).
    \item \texttt{push} (write) and \texttt{pop} (read) synchronous to the rising edge of \texttt{clk}.
    \item Data written later corresponds to earlier reads (LIFO constraint).
    \item Ignores push attempts when \texttt{full} flag is set.
    \item Ignores pop attempts when \texttt{empty} flag is set.
\end{itemize}

\subsection{Implementation (\texttt{stack.v})}
\begin{lstlisting}[language=Verilog]
`timescale 1ns / 1ps

module stack #(
    parameter WIDTH = 8,
    parameter DEPTH_LOG2 = 3
) (
    input clk, rst, push, pop,
    input [WIDTH-1:0] din,
    output reg [WIDTH-1:0] dout,
    output reg full, empty
);
    localparam MAX_DEPTH = 1 << DEPTH_LOG2;
    reg [WIDTH-1:0] mem [0:MAX_DEPTH-1];
    reg [DEPTH_LOG2:0] sp;

    always @(posedge clk or posedge rst) begin
        if (rst) begin
            sp <= 0; full <= 0; empty <= 1; dout <= 0;
        end else begin
            if (push && !full) begin
                mem[sp] <= din;
                sp <= sp + 1;
                empty <= 0;
                if (sp + 1 == MAX_DEPTH) full <= 1;
            end else if (pop && !empty) begin
                dout <= mem[sp - 1];
                sp <= sp - 1;
                full <= 0;
                if (sp - 1 == 0) empty <= 1;
            end
        end
    end
endmodule
\end{lstlisting}

\subsection{Testbench (\texttt{stack\_tb.v})}
\begin{lstlisting}[language=Verilog]
`timescale 1ns / 1ps

module stack_tb;
    parameter WIDTH = 8;
    parameter DEPTH_LOG2 = 2; // Testing depth 4

    reg clk, rst, push, pop;
    reg [WIDTH-1:0] din;
    wire [WIDTH-1:0] dout;
    wire full, empty;

    stack #(.WIDTH(WIDTH), .DEPTH_LOG2(DEPTH_LOG2)) uut (
        .clk(clk), .rst(rst), .push(push), .pop(pop), 
        .din(din), .dout(dout), .full(full), .empty(empty)
    );

    initial forever #5 clk = ~clk;

    initial begin
        rst = 1; push = 0; pop = 0; din = 0;
        #15; rst = 0; #10;

        $display("Pushing elements...");
        push_data(8'h11); push_data(8'h22); push_data(8'h33); push_data(8'h44);

        $display("Trying to push when full...");
        push_data(8'h55);
        if (full) $display("Stack is full, correctly blocked push.");

        $display("Popping elements...");
        pop_data(); pop_data(); pop_data(); pop_data();

        $display("Trying to pop when empty...");
        pop_data();
        if (empty) $display("Stack is empty, correctly blocked pop.");

        #20 $finish;
    end

    task push_data(input [WIDTH-1:0] val);
    begin
        @(posedge clk); push = 1; din = val;
        @(posedge clk); push = 0;
        $display("Pushed %h | full=%b empty=%b", val, full, empty);
    end
    endtask

    task pop_data();
    begin
        @(posedge clk); pop = 1;
        @(posedge clk); pop = 0;
        $display("Popped %h | full=%b empty=%b", dout, full, empty);
    end
    endtask
endmodule
\end{lstlisting}

\subsection{Simulation Results}
\begin{lstlisting}[language=bash]
Pushing elements...
Pushed 11 | full=0 empty=0
Pushed 22 | full=0 empty=0
Pushed 33 | full=0 empty=0
Pushed 44 | full=1 empty=0
Trying to push when full...
Pushed 55 | full=1 empty=0
Stack is full, correctly blocked push.
Popping elements...
Popped 44 | full=0 empty=0
Popped 33 | full=0 empty=0
Popped 22 | full=0 empty=0
Popped 11 | full=0 empty=1
Trying to pop when empty...
Popped 11 | full=0 empty=1
Stack is empty, correctly blocked pop.
\end{lstlisting}

\end{document}
